\section{\texorpdfstring{Derived Categories:\allowbreak{} localization and the first kernel lemmas}{Derived  localization and the first kernel lemmas}}\label{proofs-1149-1843-derived-categories-localization-and-the-first-kernel-lemmas}

Authority and comparison identities are those of \texttt{INPUT\_\allowbreak{}BINDING.\allowbreak{}json}. All unqualified Derived Categories line numbers refer to the Stacks Project authors\textquotesingle{} frozen source at \texttt{a04446e5\allowbreak{}7ec1fbc2\allowbreak{}52a871af\allowbreak{}cec7752f\allowbreak{}b2807b14}. Reading now reaches original line 1860; the boundedness definition beginning at 1845 continues past that point and is not included in this completed review interval. These are complete editorial arguments for the listed decisions,\allowbreak{} not a replacement edition of the chapter. No novelty claim is made.

\subsection{1. The localization axioms and the earlier factorization repair}\label{proofs-1149-1843-the-localization-axioms-and-the-earlier-factorization-repair}

Keep the source MS5 condition \(f\in S\iff f[1]\in S\),\allowbreak{} its MS6 condition for the third map of triangles,\allowbreak{} and both left and right fraction axioms. MS6 and identity denominators imply that every isomorphism lies in \(S\):\allowbreak{} apply MS6 to the two split triangles \((0,X,X,0,1,0)\),\allowbreak{} \((0,X,Y,0,f,0)\). The third map is forced to be \(f\). Composition with isomorphisms therefore preserves \(S\),\allowbreak{} by MS1,\allowbreak{} and comparison isomorphisms between iterated shifts do not alter membership in \(S\). In particular,\allowbreak{} for \(s\in S\),\allowbreak{} its shift \(s[-1][1]\) is conjugate to \(s\) by the specified natural isomorphisms,\allowbreak{} so \(s[-1][1]\in S\) and MS5 gives \(s[-1]\in S\). Induction yields stability under every integral shift.

For LMS2,\allowbreak{} start with \(f:X\to Y\),\allowbreak{} \(t:X\to X'\) in \(S\),\allowbreak{} and a distinguished \((X,Y,Z,f,g,h)\). Complete \(t[1]h:Z\to X'[1]\),\allowbreak{} rotating as necessary,\allowbreak{} to a triangle \((X',Y',Z,f',g',t[1]h)\). Twice rotate both triangles. Their first two vertical maps are \(1_Z,t[1]\),\allowbreak{} which belong to \(S\); MS6 gives a third map \(s'[1]\in S\). Full faithfulness of the shift and MS5 give \(s':Y\to Y'\) in \(S\) with \(s'f=f't\). The reversed construction proves RMS2. This supplies the exact shift and denominator argument at 1187–\allowbreak{}1253.

For an exact functor \((F,\xi)\),\allowbreak{} the class of arrows inverted by \(F\) satisfies MS1. It satisfies the source\textquotesingle s MS4:\allowbreak{} if \(F(fg)\) and \(F(gh)\) are invertible,\allowbreak{} then \(F(g)\) has a left inverse and a right inverse,\allowbreak{} which coincide. MS5 follows from \[
F(f[1])=\xi_Y^{-1}F(f)[1]\xi_X.
\] For MS6 choose a third map using TR3,\allowbreak{} apply the exact functor with its actual translation comparison,\allowbreak{} and use the two-isomorphisms criterion on the image triangles. The preceding paragraph gives MS2. For LMS3 retain \(a=f-g\),\allowbreak{} a triangle \((Z,X,Q,t,d,h)\),\allowbreak{} and the factorization \(a=id\) supplied by contravariant representable exactness. Completing \(i:Q\to Y\) gives \((Q,Y,W,i,j,k)\). Since \(F(t)\) is invertible,\allowbreak{} \(F(Q)=0\). The image of the latter triangle shows \(F(j)\) invertible. Thus \(j\in S\) and \(ja=jid=0\),\allowbreak{} as required. Reverse all arrows for RMS3.

The same reasoning applies to the arrows inverted by every \(H^n\) for a homological functor \(H\). Shift comparisons identify the condition at \(f[1]\) with the condition at all consecutive \(H^n(f)\). The long exact sequences give MS6 by the five lemma. For LMS3 retain the source notation at 1334–\allowbreak{}1355:\allowbreak{} \[
f,g:X\to Y,\quad a=f-g,\quad t:Z\to X,\quad
(Z,X,Q,t,u,v),\quad i:Q\to Y,\quad (Q,Y,W,i,j,k).
\] Since \(at=0\),\allowbreak{} exactness gives \textbf{\(iu=a\)}. For every integer \(n\),\allowbreak{} the exact segment \[
H^n(Z)\xrightarrow{H^n(t)}H^n(X)\longrightarrow H^n(Q)
\longrightarrow H^{n+1}(Z)\xrightarrow{H^{n+1}(t)}H^{n+1}(X)
\] has both displayed \(t\)-maps invertible,\allowbreak{} so \(H^n(Q)=0\). The long exact sequence for the second triangle now makes each \(H^n(j)\) invertible. Finally \(ja=jiu=0\).

This proves the semantic match of \textbf{OCC-01492 /\allowbreak{} P02-E0020} with \textbf{MC-STK-ERR-0812},\allowbreak{} already present in the current chapter. Its historical proof annotation calls \(g\) undefined. That explanation is inaccurate:\allowbreak{} \(g:X\to Y\) was declared at 1334. It is the wrong map for this composition; \(i\) has domain \(Q\),\allowbreak{} whereas \(u:X\to Q\) is the triangle arrow giving the factorization. The exact prior edit is correct and preserved. This note corrects the explanation without rewriting the immutable earlier receipt or creating a duplicate edit.

\subsection{2. The localization triangle morphism and two earlier copyedits}\label{proofs-1149-1843-the-localization-triangle-morphism-and-two-earlier-copyedits}

For the construction at 1394–\allowbreak{}1442 keep the original common-target diagram and all of its arrows:\allowbreak{} \[
f:X\to Y,\quad f':X'\to Y',\quad f'':X''\to Y'',\quad
k:X\to X'',\quad l:Y\to Y'',\quad s:X'\to X'',\quad t:Y'\to Y'',
\] where \(s,t\in S\),\allowbreak{} \(f''k=lf\),\allowbreak{} \(f''s=tf'\),\allowbreak{} \(a=Q(s)^{-1}Q(k)\),\allowbreak{} and \(b=Q(t)^{-1}Q(l)\). TR3 supplies \((k,l,m)\) to a completion \((X'',Y'',Z'',f'',g'',h'')\). MS6 supplies \((s,t,r)\) from the primed triangle,\allowbreak{} with \(r\in S\). Consequently \[
mg=g''l,\qquad h''m=k[1]h,
\] and \[
rg'=g''t,\qquad h''r=s[1]h'.
\] The first two arrows\textquotesingle{} equation is already \(f''k=lf\). The middle correction \textbf{MC-STK-ERR-0813} labels the fourth vertical arrow \(k[1]:X[1]\to X''[1]\),\allowbreak{} exactly as required; \(g[1]\) has the wrong source and target. This is the semantic match for \textbf{OCC-01493 /\allowbreak{} P02-E0021}.

Put \(c=Q(r)^{-1}Q(m)\). Applying \(Q\) to the displayed identities gives \[
cQ(g)=Q(r)^{-1}Q(g'')Q(l)=Q(g')Q(t)^{-1}Q(l)=Q(g')b.
\] Similarly,\allowbreak{} since the localized shift has identity comparison with \(Q\),\allowbreak{} \[
Q(h')c=Q(s[1])^{-1}Q(h'')Q(m)
=Q(s[1])^{-1}Q(k[1])Q(h)=a[1]Q(h).
\] Along with \(bQ(f)=Q(f')a\),\allowbreak{} these are all three equations of the localized triangle morphism. No map has been identified merely from a matching numerical locus.

\textbf{MC-STK-ERR-0814} at 1445 correctly replaces ``lemma'' by “proposition”:\allowbreak{} the enclosing environment is the proposition beginning at 1359. No received physical report in this interval corresponds to that older copyedit,\allowbreak{} but the unit itself has been reconciled.

\textbf{OCC-01494 /\allowbreak{} P02-E0022} matches \textbf{MC-STK-ERR-0815} at 1521:\allowbreak{} the duplicated ``is'' is already removed. Here is the required mathematical context. For the full pre-triangulated subcategory \(\mathcal D'\),\allowbreak{} the restricted denominators \(S'\) are the arrows inverted by \(Q|_{\mathcal D'}\),\allowbreak{} because \(S\) is saturated. Section 1 therefore gives their multiplicative-system and triangle compatibility properties. Every object \(X\) receives a denominator \(X'\to X\) from \(\mathcal D'\),\allowbreak{} giving essential surjectivity of the localized inclusion. For a roof \(A\leftarrow X\to B\),\allowbreak{} with \(A,B\in\mathcal D'\),\allowbreak{} choose a denominator \(X'\to X\) with \(X'\in\mathcal D'\). Composing gives an equivalent roof wholly inside the full subcategory,\allowbreak{} proving fullness. If two such roofs are equal in the ambient localization,\allowbreak{} choose a common right refinement on which their numerator maps agree. Apply the same density assumption to its source; the resulting refinement lies in \(\mathcal D'\),\allowbreak{} its denominators lie in \(S'\),\allowbreak{} and the numerators remain equal. This proves faithfulness as well as the cofinality implicit in the source\textquotesingle s colimit argument.

\subsection{\texorpdfstring{3. DERIVED-NEW-005:\allowbreak{} fix the comparison in the uniqueness assertion}{3.  fix the comparison in the uniqueness assertion}}\label{proofs-1149-1843-derived-new-005-fix-the-comparison-in-the-uniqueness-assertion}

At 1363–\allowbreak{}1365 the uniqueness assertion requires that the localization functor\textquotesingle s given translation comparison be the identity under \([1]Q=Q[1]\). An exact functor in the chapter\textquotesingle s definition includes that comparison; equality of the two composite functors does not itself force every natural automorphism between them to be the identity. The intended structure in the proof has the identity comparison.

With it fixed,\allowbreak{} the proof is exact. MS5 lets the shift act on a left fraction by \[
\bigl(Q(s)^{-1}Q(f)\bigr)[1]=Q(s[1])^{-1}Q(f[1]).
\] The localization universal property makes this well-defined and unique; the inverse shift descends in the same way. Canonical distinguished triangles are the isomorphic images of distinguished triangles of \(\mathcal D\). If a second pre-triangulated structure on this fixed localized additive category and shift makes \((Q,1)\) exact,\allowbreak{} it contains all these canonical triangles. Any distinguished triangle of the second structure has some first arrow \(a\). A canonical completion of \(a\) exists by the source\textquotesingle s TR1 construction. Both triangles now belong to the second structure and share that first arrow,\allowbreak{} so TR3 and representable exactness give an isomorphism between them with identity first two components. Isomorphism closure of the canonical class puts the original triangle in the canonical class. Thus the two classes agree. This proves uniqueness with the exact comparison specified,\allowbreak{} rather than adding an unproved uniqueness claim.

To see the limitation of leaving the comparison free,\allowbreak{} let \((\mathcal D,[1],\mathcal T)\) be a triangulated category and define \[
\mathcal T^-=
\{(X,Y,Z,f,g,-h):(X,Y,Z,f,g,h)\in\mathcal T\}.
\] This is another triangulation on the same additive category and shift. TR1 follows directly. For TR2,\allowbreak{} the rotation of \((X,Y,Z,f,g,-h)\) is \((Y,Z,X[1],g,-h,-f[1])\). It is in \(\mathcal T^-\) exactly when \((Y,Z,X[1],g,-h,f[1])\) is in \(\mathcal T\). The latter is isomorphic via \((1,1,-1)\) to \((Y,Z,X[1],g,h,-f[1])\); TR2 for \(\mathcal T\) gives the required equivalence in both directions. For TR3,\allowbreak{} negating both boundary maps leaves the third commuting-square equation equivalent to the original one,\allowbreak{} so the same third map works.

For TR4 start with three \(\mathcal T^-\)-triangles with boundary maps \(d_1,d_2,d_3\). Their corresponding \(\mathcal T\)-triangles have boundaries \(-d_1,-d_2,-d_3\). TR4 in \(\mathcal T\) provides \(a,b\) with connecting arrow \(-p_1[1]d_3\). Negating that arrow gives the required \(\mathcal T^-\)-triangle with connecting arrow \(p_1[1]d_3\). The other two triangle-morphism conditions still hold because both boundaries in each equation were negated. This proves all four axioms.

Now \((\operatorname{id},-1_{[1]}):
(\mathcal D,\mathcal T)\to(\mathcal D,\mathcal T^-)\) is an exact functor and commutes strictly with the shift as a functor. Take \(S\) to be all isomorphisms,\allowbreak{} so that the identity is a model of the localization functor. Both structures satisfy an assertion requiring only an unspecified exact structure on that functor. They need not coincide. In the cone example of \texttt{PROOFS\_\allowbreak{}0001\_\allowbreak{}1140.\allowbreak{}md},\allowbreak{} §3,\allowbreak{} the triangle \((\mathbb Z[0],\mathbb Z[0],C(3),3,i,-p)\) belongs to \(\mathcal T\),\allowbreak{} whereas \((\mathbb Z[0],\mathbb Z[0],C(3),3,i,p)\) does not:\allowbreak{} comparison with identity first two maps would require an integer endomorphism class \(m\) of the cone to satisfy both \(m\equiv1\pmod3\) and \(m\equiv-1\pmod3\). The unshifted homotopy differential is \(3\),\allowbreak{} so homotopies change components by \(3t\),\allowbreak{} giving exactly those congruences. This proves distinctness. Transport along the canonical equivalence gives the same example on the particular fraction model of \(S^{-1}\mathcal D\).

This is a concrete consequence of the preceding sign finding and the chapter\textquotesingle s own definition of an exact functor. The repair makes explicit the intended comparison; it does not change the constructed localization.

The four external current-chapter citations to this proposition were checked at \texttt{adequate.\allowbreak{}tex:\allowbreak{}2202–\allowbreak{}2227},\allowbreak{} \texttt{dga.\allowbreak{}tex:\allowbreak{}2568–\allowbreak{}2590},\allowbreak{} \texttt{perfect.\allowbreak{}tex:\allowbreak{}9810–\allowbreak{}9840},\allowbreak{} and \texttt{sdga.\allowbreak{}tex:\allowbreak{}4187–\allowbreak{}4212}. They use existence of the canonical localization or its exact universal factorization,\allowbreak{} which uses precisely \((Q,1)\). None needs uniqueness with a free translation comparison. The Perfect Complexes passage also uses the reflection lemma; its exact functor retains its own specified comparison as in DERIVED-NEW-002. No change to these receiving passages is required. Later internal citations are still checked in chapter order.

The exact structure on the factorization \(F'\) at 1483–\allowbreak{}1488 also has a definite receiving map. Since localization has the same objects,\allowbreak{} put \(\xi'_{Q(X)}=\xi_X\). Naturality on \(Q(f)\) is the original naturality,\allowbreak{} and naturality on inverses \(Q(s)^{-1}\) follows by multiplying that same commuting equation by inverses. It therefore holds on every fraction. The images of canonical distinguished triangles are distinguished,\allowbreak{} so \((F',\xi')\) is exact. This proves the relevant part of the omitted universal-property argument with the translation maps exposed.

\subsection{4. Kernel witnesses and filtered triangle denominators}\label{proofs-1149-1843-kernel-witnesses-and-filtered-triangle-denominators}

There is no new error at 1574–\allowbreak{}1577. Given the zero denominator \(0:Z\to Z'\),\allowbreak{} the source\textquotesingle s split triangle \[
(Z'[-1],Z'[-1]\oplus Z,Z,\iota,\pi,0)
\] rotates \textbf{backwards} to \[
(Z[-1],Z'[-1],Z'[-1]\oplus Z,-0[-1],\iota,\pi).
\] The first arrow is in \(S\) by MS5 and closure under comparison isomorphisms; its minus sign does not change a zero map. The third object is isomorphic,\allowbreak{} by the actual biproduct interchange map,\allowbreak{} to \(Z\oplus Z'[-1]\). This is exactly condition (4)\allowbreak{},\allowbreak{} using \(Z'[-1]\) as the complementary object. A forward rotation alone would expose a shifted copy of \(Z\); the inverse rotation proves the stated unshifted result. Conversely,\allowbreak{} if \(f\in S\) and its triangle cone is \(Z\oplus Z'\),\allowbreak{} exactness of \(Q\) gives \(Q(Z)\oplus Q(Z')=0\). Each summand\textquotesingle s identity factors through the zero object,\allowbreak{} so each summand is zero. For saturated \(S\),\allowbreak{} the two zero-object arrows are denominators exactly when their localized images are invertible,\allowbreak{} giving conditions (5)\allowbreak{} and (6)\allowbreak{}; the triangle \((0,Z,Z,0,1,0)\) gives condition (7)\allowbreak{}.

For the category \(\mathcal I\) at 1595–\allowbreak{}1735,\allowbreak{} an object is the actual denominator triangle map \(s:T\to T'\),\allowbreak{} with all three components in \(S\); a morphism is a triangle map commuting with these structure maps. Every component of such a morphism is in \(S\):\allowbreak{} its image under \(Q\) is the quotient of two invertible structure maps,\allowbreak{} and \(S\) is saturated. Thus subsequent denominator completions used in the proof really have the required membership.

MS2 followed by MS6 completes any prescribed denominator \(X\to X'\) to an object of \(\mathcal I\). Rotation gives the same assertion for either other component. Applied to a target triangle,\allowbreak{} this also completes any subsequent denominator on one component to a triangle map with all three components in \(S\). The latter assertion permits successive replacements of a target triangle; it does not prescribe its two other new objects in advance.

For two objects \(s_1:T\to T_1\),\allowbreak{} \(s_2:T\to T_2\),\allowbreak{} filteredness of each component category \(X/S,Y/S,Z/S\) and these completions give a common later target and three maps \(a,b,c\) from \(T_1\) to that target,\allowbreak{} commuting with the three maps out of \(T\). Each remaining triangle-square defect becomes zero after \(Q\):\allowbreak{} both triples there are the quotient of the same two structure isomorphisms. The fraction equality criterion gives a denominator out of the target component that kills one defect. Complete it to a map of target triangles,\allowbreak{} postcompose the triple,\allowbreak{} and do this for each of its other two components,\allowbreak{} using rotations when the defect lands in a shifted component. Once a square commutes,\allowbreak{} postcomposition by a triangle map preserves its commutativity. After these three steps all squares commute,\allowbreak{} giving the required common object of \(\mathcal I\).

For two parallel arrows in \(\mathcal I\),\allowbreak{} their first components agree after composition with the first denominator of their source. LMS3 gives a target denominator making those components equal. Complete it to a map of triangles and then do the same for the second and third components. Equalities already obtained survive composition,\allowbreak{} so the resulting two triples are equal. The identity of \(T\) supplies a nonempty object. These arguments verify all three filteredness axioms.

Finally let \(u:X\to U\) be an object of \(X/S\),\allowbreak{} with two arrows to the first components of two objects of \(\mathcal I\). Choose a common target in \(\mathcal I\). The two resulting arrows from \(U\) to its first component can be equalized in the filtered category \(X/S\). Complete that equalizing denominator to one more target triangle map. This supplies a common later object with equal induced arrows out of \(U\),\allowbreak{} giving a connected comma category. The earlier completion gives its nonemptiness. This proves cofinality in the source\textquotesingle s exact definition; rotations prove the other two cofinality assertions. No claim that a surjective-on-objects functor alone is cofinal is used.

\subsection{\texorpdfstring{5. DERIVED-NEW-006:\allowbreak{} propagate the image-triangle comparison}{5.  propagate the image-triangle comparison}}\label{proofs-1149-1843-derived-new-006-propagate-the-image-triangle-comparison}

The same typed-boundary repair as DERIVED-NEW-002 applies at 1787 and 1803. Name the exact functor \((F,\xi)\),\allowbreak{} and write the image triangle with boundary \(\xi_XF(h)\). The kernel conclusion stays unchanged:\allowbreak{} naturality and invertibility of \(\xi\) carry zero objects through both shifts; the image triangle with its first two objects zero has third object zero by representable exactness. Fullness of the kernel subcategory and TR3 supply all of its triangle maps; the cone criterion then makes it pre-triangulated,\allowbreak{} and TR4 restricts when available in the ambient category. If an object \(X\oplus Y\) maps to zero,\allowbreak{} additivity of \(F\) gives \(F(X)\oplus F(Y)=0\); projection and inclusion show the identities of both summands are zero. Hence the kernel is saturated and strictly full. The same argument applied to every \(H^n\) proves the following homological-kernel lemma at 1812–\allowbreak{}1843,\allowbreak{} using its long exact sequence instead of an image triangle.

The new operations in this interval are one comparison qualification in the localization proposition and two explicit-map operations in the kernel lemma. Four earlier units remain unchanged; three further physical reports are accounted for by those earlier corrections.
